\section{Methodology}

\subsection{Candidate paper generation}
\label{sec:candidate-generation}

Papers enter this study through a single, reusable candidate-generation pipeline, run
once per side of the comparison, prior to the individual manual review documented in
Section~\ref{sec:results-detail}.

\subsubsection{Category-based corpus construction}
\label{sec:category-corpus-construction}

Every candidate paper on both sides is retrieved by a single, comprehensive pull
against arXiv's own subject categories, with no keyword or search-term filtering of
any kind. Each side is queried against a fixed set of ``home'' categories --- the
categories where that side's research actually concentrates --- over a 2021-present
submission-date window:

\begin{itemize}[leftmargin=1.5em]
    \item \textbf{Quantum Computing}: \texttt{quant-ph}, \texttt{cond-mat.supr-con},
    \texttt{cond-mat.mes-hall}.
    \item \textbf{Classical computing}: \texttt{cs.DS}, \texttt{cs.CC},
    \texttt{cs.LG}, \texttt{cs.NA}, \texttt{cs.IT}, \texttt{cs.CR}, \texttt{cs.AR},
    \texttt{cs.PL}, \texttt{cs.DC}, \texttt{cs.ET}, \texttt{eess.SY},
    \texttt{eess.SP}, \texttt{physics.app-ph}.
\end{itemize}

These are each side's own home categories specifically, not the full set of
categories either side's subject matter can ever touch --- e.g.\ the quantum side
does not include \texttt{cs.LG}, since that would import the entire classical
machine-learning literature into the quantum-side corpus. Because no keyword filter
is applied within these categories, retrieval captures every paper in scope
regardless of how its authors chose to phrase their own topic.

\paragraph{Deep-pagination limit and fix.} arXiv's search API returns a hard error
for any query offset at or beyond approximately 10{,}000 results, confirmed directly
(an identical offset succeeds at 9{,}950 and fails at 10{,}000, 10{,}050, 20{,}000,
and 30{,}000, independent of retry count or backoff) --- a structural ceiling in
arXiv's own search backend, not a transient fault. Since both sides' full category
pulls exceed this ceiling by a wide margin, retrieval recursively splits the
submission-date range: a candidate date window is first probed for its total result
count; a window safely under the ceiling is paginated directly, while a window still
too large is split in half by date and each half is probed and paginated
independently, continuing until every actual paginated request stays comfortably
clear of the limit. This adapts automatically to uneven publication volume across
time (recent months have far more submissions than 2021 months) rather than assuming
a fixed window size that might still overflow for a heavy period.

\paragraph{Result.} This procedure produces a complete candidate corpus of
\textbf{109{,}321 quantum papers and 345{,}107 classical-computing papers}
(454{,}428 papers total), verified free of duplicate arXiv identifiers on both sides.
Table~\ref{tab:corpus-stats} reports these totals alongside current classification
status.

\subsubsection{Per-paper extraction and link discovery}
\label{sec:candidate-list-gen}

Every paper in a domain's corpus goes through the same deterministic extraction pass;
nothing is pre-filtered or skipped at this stage, and nothing here decides which
papers get read later.

\paragraph{Extraction.} For each paper, the PDF is downloaded from arXiv, its full
plain text is extracted (\texttt{pdftotext}), and its embedded hyperlink annotations
are extracted separately (\texttt{pdftohtml}, run with image extraction disabled ---
the corpus needs only text, and image extraction is unnecessary I/O and storage
overhead at this scale). The PDF itself is deleted immediately afterward; only the
extracted text and link list are retained, since keeping the full-corpus PDFs is not
disk-feasible at this scale. Before pattern matching, the extracted text is
\emph{flattened}: hyphenated words and URLs split across a PDF line break are
rejoined, so a URL or DOI broken across a line wrap (e.g.\ \texttt{https:} at the end
of one line, \texttt{//doi.org/...} at the start of the next) is still read as one
continuous string.

\paragraph{Link discovery.} The same pass that extracts a paper's text also finds
every plausible data/code link in it, from three independent sources: (i) PDF
hyperlink annotations and plain-text URLs matching \textasciitilde20 known data/code
repository domains (GitHub, Zenodo, Figshare, OSF, Dryad, institutional
research-data repositories, etc.); (ii) generic DOI and URL patterns scanned
independently of nearby text, to catch a repository link with no availability phrase
next to it; and (iii) proximity to \textasciitilde30 data/code-availability keyword
phrases (e.g.\ \emph{data availab-}, \emph{deposited at}, \emph{openly available}),
used only to flag where in the text a link or claim sits. Every paper's discovered
links proceed to the next stage (Section~\ref{sec:extraction-classification}), where
each one is independently checked --- not only for whether it resolves, but, for the
platforms that matter most, whether it actually has data behind it. Distinguishing a
genuine self-citation (e.g.\ an author linking their own prior repository) from a
citation to an unrelated third-party tool is left to the verification and
classification steps that follow, not decided here.

\paragraph{Corpus size and classification status at time of writing.}
\label{sec:corpus-size}

Table~\ref{tab:corpus-stats} reports the current corpus size on each side and how
many papers have completed the two-axis classification described in
Section~\ref{sec:extraction-classification}.

\begin{table}[htbp]
\centering
\caption{Corpus size and classification status by domain}
\label{tab:corpus-stats}
\begin{tabular}{@{}l r r@{}}
\toprule
Domain & Corpus size & Classified so far \\
\midrule
Quantum & 109{,}321 & 3{,}472 \\
Classical computing (CS) & 345{,}107 & 0 \\
\bottomrule
\end{tabular}

\vspace{2mm}
\raggedright\footnotesize
Full-text extraction, link discovery, and classification
(Sections~\ref{sec:candidate-list-gen} and~\ref{sec:extraction-classification}) are
still in progress against the corpus reported above; per-domain results in this
report will be updated as classification proceeds.

\end{table}
